% !PRINTED 245
\weberpage{245}

\weq{3}{a_1 < 0, \quad a_2 > 0, \quad a_3 < 0;}
diese Bedingungen sind nothwendig dafür, dass alle drei Wurzeln
positiv sind. Sie sind aber auch hinreichend; denn wenn eine
oder drei Wurzeln negativ sind, so ist $a_3 > 0$. Sind aber zwei
Wurzeln, etwa $\beta, \gamma$, negativ, so ist entweder
\[
a_1 \geqq 0 \quad \text{oder} \quad \alpha > -(\beta + \gamma).
\]
In letzterem Falle aber folgt
\[
a_2 < \beta\gamma - (\beta + \gamma)^2 = -\beta^2 - \beta\gamma - \gamma^2,
\]
also $a_2$ negativ.

Ist endlich eine Wurzel gleich 0, so muss nothwendig $a_3$
verschwinden.

Indem man $x$ durch $-x$ ersetzt, schliesst man, dass für drei
negative Wurzeln die nothwendige und hinreichende Bedingung die
ist
\weq{4}{a_1 > 0, \quad a_2 > 0, \quad a_3 > 0.}
Wir erhalten daher unter Voraussetzung einer positiven
Discriminante folgende Tabelle:
\[
\begin{array}{ccc|c}
a_1 & a_2 & a_3 & \\ \hline
- & + & - & 3\\[2pt]
- & - & - & \\
+ & + & - & \raisebox{1.2em}{$\big\}1$}\\
+ & - & - & \\ \hline
+ & + & + & 0\\[2pt]
- & - & + & \\
- & + & + & \raisebox{1.2em}{$\big\}2$}\\
+ & - & + &
\end{array}
\]
wo in der letzten Columne die Zahl der positiven Wurzeln angegeben
ist. Wir können das Resultat dieser Betrachtung so aussprechen:

\gesperrt{Bei positiver Discriminante ist die Anzahl der
positiven Wurzeln der cubischen Gleichung gleich der Anzahl der
Zeichenwechsel in der Reihe}
\[
1, \; a_1, \; a_2, \; a_3,
\]
wenn wir unter einem Zeichenwechsel die Aufeinanderfolge einer
positiven und einer negativen oder einer negativen und einer
positiven Grösse verstehen.
